\item La figura P34.21 muestra un rayo de luz incidente en una serie de placas que tienen diferentes índices de refracción, donde $n_1 < n_2 < n_3 < n_4$. Observe que la trayectoria del rayo se inclina hacia la normal. Si la variación en $n$ fuera continua, la ruta podría formar una curva suave. Use esta idea y un diagrama de rayos para explicar por qué se puede ver el Sol al atardecer después de que físicamente ha caído por debajo del horizonte.